\chapter{Gestion economica}

\section{Recursos humanos}
Aunque el proyecto es individual, se modelan roles profesionales para estimar el coste realista del trabajo: jefe de proyecto, analista de sistemas, desarrollador y tester. Esta separación permite valorar de forma más rigurosa tareas de planificación, diseño, implementación y validación.

\section{Recursos materiales directos}
El coste material principal corresponde al hardware de inferencia y servicios: servidor doméstico, GPU AMD con 16 GB de VRAM, almacenamiento NVMe y dispositivos de bajo consumo como Raspberry Pi para tareas 24/7. La memoria final deberá actualizar esta sección con el hardware realmente utilizado, precios finales y justificación de cada componente.

\section{Recursos materiales indirectos}
Se incluyen conectividad, consumo energético, espacio de trabajo y amortización de equipos. Un punto relevante del proyecto es la comparación entre un modelo de inversión inicial en infraestructura propia y un modelo de suscripciones o consumo cloud.

\section{Recursos software}
El ecosistema se apoya mayoritariamente en software libre o de código abierto: Linux, Docker, motores de inferencia local, herramientas de automatización, servicios autoalojados y LaTeX para la documentación. El coste directo de licencias se estima inicialmente en cero, aunque se documentarán posibles costes indirectos de mantenimiento, aprendizaje y tiempo de configuración.

\section{Contingencias e imprevistos}
La planificación contempla una reserva para desviaciones técnicas: incompatibilidades de drivers, saturación de memoria gráfica, fallos de hardware, cambios de arquitectura y reducción de alcance. En la memoria final se diferenciará entre contingencia genérica e imprevistos derivados de riesgos identificados.

\section{Estimacion final del presupuesto}
La planificación inicial estimó un presupuesto total de 14.608,34 euros, incluyendo costes de personal, costes genéricos, contingencia e imprevistos. Esta cifra se revisará con datos reales y se explicará como presupuesto imputado de proyecto académico, no como desembolso directo completo.

\section{Control de gestion}
El control económico se realizará al cierre de cada sprint, comparando horas estimadas y horas reales, además de registrar cambios relevantes de alcance. Las desviaciones negativas se asociarán a riesgos concretos cuando sea posible para justificar técnicamente el impacto.
